\chapter{结论}

基于溯源图的威胁检测已经从规则匹配和统计异常发展到图神经网络、自监督学习、在线细粒度检测、LLM辅助调查和EDR工程化落地。总体趋势可以概括为三点：第一，检测粒度从图级异常走向节点、路径和子图级归因；第二，系统目标从追求单一准确率走向同时优化实时性、误报、解释性和人工工作量；第三，研究边界从检测模型扩展到数据集建设、日志到图的自动构建、语义补全和企业部署。

然而，该领域距离成熟落地仍有距离。现有benchmark过度集中于DARPA系列，攻击技术和业务场景更新不足；许多方法在统一评测框架下表现并不稳定；机器学习方法面临概念漂移、数据投毒和可解释性压力；LLM方法虽然改善语义解释，却必须解决事实约束、成本和隐私问题。未来研究应围绕可复现评测、跨系统泛化、高质量归因和人机协同调查展开，使溯源图真正成为连接底层审计日志和高层安全决策的可靠桥梁。
